\documentclass[10pt]{article}
\usepackage[T1]{fontenc}
\usepackage[utf8]{inputenc}
\usepackage[french]{babel}
\usepackage{amsmath,amssymb,mathrsfs,mathpazo}
\usepackage{geometry}
\geometry{a4paper,margin=2cm}
\usepackage{graphicx,array,tabularx,tabularray,colortbl,multirow,multicol,booktabs,enumitem,pifont,ulem,fancybox,diagbox,cancel}
\usepackage{pstricks,pstricks-add,pst-plot,pst-tree,pst-node,pst-text,pst-func,pst-math,pst-bspline,pst-3dplot,pst-eucl,pst-solides3d}
\let\pscalculate\undefined
\usepackage{pst-bezier}
\usepackage{tikz,pgfplots,tkz-tab}
\usetikzlibrary{arrows,arrows.meta,patterns,calc,positioning,trees,decorations.markings,shapes.geometric,shapes.misc,shapes.arrows,intersections,decorations.pathmorphing,decorations.pathreplacing,backgrounds,shadows,babel}
\pgfplotsset{compat=1.18}
\usepgfplotslibrary{fillbetween}
\usepackage[np]{numprint}
\usepackage{listings,makecell,pas-tableur,twoopt}
\let\barre\undefined
\usepackage[pstricks]{tablvar}
\definecolor{darkgreen}{RGB}{0,100,0}
\providecommand{\red}{\color{red}}
\providecommand{\blue}{\color{blue}}
\providecommand{\green}{\color{green}}
\providecommand{\black}{\color{black}}
\providecommand{\white}{\color{white}}
\providecommand{\gray}{\color{gray}}
\providecommand{\cyan}{\color{cyan}}
\providecommand{\magenta}{\color{magenta}}
\providecommand{\yellow}{\color{yellow}}
\providecommand{\violet}{\color{violet}}
\providecommand{\orange}{\color{orange}}
\providecommand{\euro}{\texteuro}
\providecommand{\EUR}{\texteuro}
\providecommand{\ieme}{\textsuperscript{e}}
\providecommand{\degres}{\ensuremath{^\circ}}
\providecommand{\degree}{\ensuremath{^\circ}}
\providecommand{\vv}[1]{\overrightarrow{#1}}
\providecommand{\Coord}[1]{#1}
\providecommand{\decosix}{}
\providecommand{\starredbullet}{\ensuremath{\bullet}}
\pagestyle{empty}
\setlength{\parindent}{0pt}
\usepackage{adjustbox}\DeclareUnicodeCharacter{2236}{:}\DeclareUnicodeCharacter{2190}{\ensuremath{\leftarrow}}
\providecommand{\R}{\mathbb{R}}
\providecommand{\N}{\mathbb{N}}
\providecommand{\Z}{\mathbb{Z}}
\providecommand{\Q}{\mathbb{Q}}
\providecommand{\D}{\mathbb{D}}
\providecommand{\C}{\mathbb{C}}
\providecommand{\e}{{\mathrm{e}}}
\providecommand{\vect}[1]{\overrightarrow{#1}}
\providecommand{\vectt}[1]{\overrightarrow{\mathrm{#1}}}
\providecommand{\vvt}[1]{\overrightarrow{\mathrm{#1}}}
\providecommand{\np}[2][]{{#2~#1}}
\providecommand{\vv}[1]{{\overrightarrow{#1}}}
\renewcommand{\d}{{\mathrm{d}}}
\providecommand{\degres}{{^\circ}}
\providecommand{\degre}{{^\circ}}
\providecommand{\blue}{\color{blue}}
\providecommand{\red}{\color{red}}
\providecommand{\enskip}{\quad}
\providecommand{\llbracket}{\lbrack\!\lbrack}
\providecommand{\rrbracket}{\rbrack\!\rbrack}
\providecommand{\euro}{{\text{€}}}
\providecommand{\no}{{\text{n°}}}
\providecommand{\interval}[2]{{[#1;#2]}}
\providecommand{\intervFO}[2]{{[#1;#2[}}
\providecommand{\dotfill}{\ldots}

\providecommand{\ds}{\displaystyle}
\providecommand{\textding}[1]{#1}
\providecommand{\Oij}{\(\mathrm{O};\vec{\imath},\vec{\jmath}\)}
\providecommand{\Oijk}{\(\mathrm{O};\vec{\imath},\vec{\jmath},\vec{k}\)}
\providecommand{\Ouv}{\(\mathrm{O};\vec{u},\vec{v}\)}
\providecommand{\intervFF}[2]{\ensuremath{[#1;#2]}}
\renewcommand{\euro}{\ifmmode\text{\texteuro}\else\texteuro\fi}
\DeclareUnicodeCharacter{00B0}{\ifmmode^{\circ}\else\textdegree\fi}
\DeclareUnicodeCharacter{20AC}{\ifmmode\text{\texteuro}\else\texteuro\fi}
% Énoncés : Éducation nationale. Transcription et corrigés : APMEP, Denis Vergès et contributeurs. Découpage et classement : CoopMaths, Sébastien Lozano, Éric Elter et contributeurs. Adaptation de la présentation : Maths Explorer.
% CC BY-SA : https://coopmaths.fr/www/about/
% https://coopmaths.fr/alea/static/eam/2026/tex/eam_specialite_2026_06_etrangers_2_cor.tex
\begin{document}
\noindent\textit{APMEP / CoopMaths. CC BY-SA. Présentation adaptée par Maths Explorer.}\par\medskip
On se place dans le plan muni d'un repère orthonormé.

On considère les points $A(-1\ ;\ 5)$, $B(3\ ;\ 5)$ et $C(4\ ;\ 0)$.

\begin{enumerate}
\item \begin{enumerate}
\item On calcule les coordonnées des vecteurs $\overrightarrow{AB}$ et $\overrightarrow{AC}$.

On rappelle que, si $M(x_M\ ;\ y_M)$ et $N(x_N\ ;\ y_N)$, alors :
\[
\overrightarrow{MN}(x_N-x_M\ ;\ y_N-y_M).
\]

Ainsi :
\begin{align*}
\overrightarrow{AB}
&=(x_B-x_A\ ;\ y_B-y_A)\\
&=(3-(-1)\ ;\ 5-5)\\
&=(4\ ;\ 0).
\end{align*}

De même :
\begin{align*}
\overrightarrow{AC}
&=(x_C-x_A\ ;\ y_C-y_A)\\
&=(4-(-1)\ ;\ 0-5)\\
&=(5\ ;\ -5).
\end{align*}

Donc :
\[
\overrightarrow{AB}(4\ ;\ 0)
\qquad \text{et} \qquad
\overrightarrow{AC}(5\ ;\ -5).
\]

\item On utilise l'expression du produit scalaire avec les coordonnées.

Si $\vec u(x\ ;\ y)$ et $\vec v(x'\ ;\ y')$, alors :
\[
\vec u\cdot \vec v=xx'+yy'.
\]

Donc :
\begin{align*}
\overrightarrow{AB}\cdot\overrightarrow{AC}
&=4\times 5+0\times (-5)\\
&=20+0\\
&=20.
\end{align*}

Ainsi :
\[
\overrightarrow{AB}\cdot\overrightarrow{AC}=20.
\]
\end{enumerate}

\item \begin{enumerate}
\item On calcule la longueur $AC$ à l'aide de la formule de distance dans un repère orthonormé.

On a :
\begin{align*}
AC
&=\sqrt{(x_C-x_A)^2+(y_C-y_A)^2}\\
&=\sqrt{(4-(-1))^2+(0-5)^2}\\
&=\sqrt{5^2+(-5)^2}\\
&=\sqrt{25+25}\\
&=\sqrt{50}\\
&=\sqrt{25\times 2}\\
&=5\sqrt{2}.
\end{align*}

Donc :
\[
AC=5\sqrt{2}.
\]

\item On admet que $AB=4$.

L'expression du produit scalaire en fonction de l'angle $\widehat{BAC}$ est :
\[
\overrightarrow{AB}\cdot\overrightarrow{AC}
=AB\times AC\times \cos\left(\widehat{BAC}\right).
\]

En remplaçant $AB$ par $4$ et $AC$ par $5\sqrt{2}$, on obtient :
\[
\overrightarrow{AB}\cdot\overrightarrow{AC}
=4\times 5\sqrt{2}\times \cos\left(\widehat{BAC}\right).
\]

C'est-à-dire :
\[
\overrightarrow{AB}\cdot\overrightarrow{AC}
=20\sqrt{2}\cos\left(\widehat{BAC}\right).
\]

\item D'après la question 1.b, on sait que :
\[
\overrightarrow{AB}\cdot\overrightarrow{AC}=20.
\]

Or, d'après la question précédente :
\[
\overrightarrow{AB}\cdot\overrightarrow{AC}
=20\sqrt{2}\cos\left(\widehat{BAC}\right).
\]

On obtient donc :
\begin{align*}
20&=20\sqrt{2}\cos\left(\widehat{BAC}\right)\\
\cos\left(\widehat{BAC}\right)&=\frac{20}{20\sqrt{2}}\\
\cos\left(\widehat{BAC}\right)&=\frac{1}{\sqrt{2}}\\
\cos\left(\widehat{BAC}\right)&=\frac{\sqrt{2}}{2}.
\end{align*}

L'angle $\widehat{BAC}$ est un angle géométrique ; sa mesure appartient donc à l'intervalle $[0\ ;\ \pi]$.

Comme :
\[
\cos\left(\frac{\pi}{4}\right)=\frac{\sqrt{2}}{2},
\]
on en déduit que :
\[
\widehat{BAC}=\frac{\pi}{4}.
\]

Une mesure, en radian, de l'angle $\widehat{BAC}$ est donc :
\[
\boxed{\frac{\pi}{4}}.
\]
\end{enumerate}
\end{enumerate}

\clearpage

\end{document}
